% "Class Notes Template" by Siladittya Manna,
% https://www.overleaf.com/latex/templates/class-notes-template/jrwyqnwjbkfr
% licensed CC BY 4.0 (https://creativecommons.org/licenses/by/4.0/).
% Changes: this comment; the metadata and the placeholder text, replaced
% with a short lecture of our own.
\documentclass[11pt]{article}
\usepackage[margin=1in]{geometry}

% Mathematics and figures
\usepackage{amsmath,amssymb}
\usepackage{graphicx}
\usepackage{tikz-cd}
\usepackage{multicol}
\usepackage[most]{tcolorbox}
\usepackage{xcolor}

% Paragraph formatting
\setlength{\parindent}{0pt}
\setlength{\parskip}{1\baselineskip}

% --------------------------------------------------------------------------
% Class-note colors
% --------------------------------------------------------------------------
\definecolor{QuestionBlue}{HTML}{2563A6}
\definecolor{StatementGray}{HTML}{59636E}
\definecolor{DefinitionTeal}{HTML}{177E89}
\definecolor{TheoremPurple}{HTML}{6B4C9A}
\definecolor{PropositionBlue}{HTML}{3F5F9A}
\definecolor{ProofGray}{HTML}{6C757D}
\definecolor{AnswerGreen}{HTML}{2E7D32}
\definecolor{ExampleOrange}{HTML}{C56A16}
\definecolor{RemarkGold}{HTML}{9A7611}

% Shared appearance for all class-note boxes
\tcbset{
  classnote/.style={
    enhanced,
    breakable,
    boxrule=0.7pt,
    arc=2pt,
    left=8pt,
    right=8pt,
    top=6pt,
    bottom=6pt,
    before skip=10pt,
    after skip=10pt,
    fonttitle=\bfseries
  }
}

% --------------------------------------------------------------------------
% Reusable class-note environments
% The optional argument changes the displayed title.
% --------------------------------------------------------------------------
\newtcolorbox{question}[1][Question]{
  classnote, colback=QuestionBlue!5, colframe=QuestionBlue, title={#1}
}
\newtcolorbox{statement}[1][Statement]{
  classnote, colback=StatementGray!5, colframe=StatementGray, title={#1}
}
\newtcolorbox{definition}[1][Definition]{
  classnote, colback=DefinitionTeal!5, colframe=DefinitionTeal, title={#1}
}
\newtcbtheorem[auto counter]{theorem}{Theorem}{
  classnote, colback=TheoremPurple!5, colframe=TheoremPurple
}{thm}
\newtcbtheorem[auto counter]{proposition}{Proposition}{
  classnote, colback=PropositionBlue!5, colframe=PropositionBlue
}{prop}
\newtcolorbox{proof}[1][Proof]{
  classnote, colback=ProofGray!5, colframe=ProofGray, title={#1},
  after upper={\hfill$\square$}
}
\newtcolorbox{answer}[1][Answer]{
  classnote, colback=AnswerGreen!5, colframe=AnswerGreen, title={#1}
}
\newtcolorbox{example}[1][Example]{
  classnote, colback=ExampleOrange!5, colframe=ExampleOrange, title={#1}
}
\newtcolorbox{remark}[1][Remark]{
  classnote, colback=RemarkGold!5, colframe=RemarkGold, title={#1}
}

% --------------------------------------------------------------------------
% Metadata: edit these fields for each set of class notes.
% --------------------------------------------------------------------------
\newcommand{\notestitle}{Convexity and Jensen's Inequality}
\newcommand{\courseid}{STAT 101}
\newcommand{\instructorname}{A.~Lecturer}
\newcommand{\classdate}{October 7}
\newcommand{\documentlabel}{Classroom Queries}

% Logo files. Paths must remain relative to this TeX file.
\newcommand{\courselogo}{course.jpeg}
\newcommand{\institutelogo}{inst.png}

% --------------------------------------------------------------------------
% Class-note header
% --------------------------------------------------------------------------
\newcommand{\makeclassnotesheader}{%
\begin{tcolorbox}[
  enhanced,
  colback=QuestionBlue!3,
  colframe=QuestionBlue,
  boxrule=1.1pt,
  arc=3pt,
  left=12pt,
  right=12pt,
  top=12pt,
  bottom=10pt
]
\noindent
\begin{minipage}[c]{0.16\linewidth}
  \centering
  \IfFileExists{\courselogo}
    {\includegraphics[width=0.94\linewidth,keepaspectratio]{\courselogo}}
    {\fbox{\scriptsize Course Logo}}
\end{minipage}%
\hfill
\begin{minipage}[c]{0.62\linewidth}
  \centering
  {\small\bfseries\color{QuestionBlue}\MakeUppercase{\documentlabel}}\\[5pt]
  {\LARGE\bfseries \notestitle}
\end{minipage}%
\hfill
\begin{minipage}[c]{0.16\linewidth}
  \centering
  \IfFileExists{\institutelogo}
    {\includegraphics[width=0.94\linewidth,keepaspectratio]{\institutelogo}}
    {\fbox{\scriptsize Institute Logo}}
\end{minipage}
\vspace{5pt}
\hrule
\vspace{8pt}
\begin{tabular*}{\linewidth}{@{\extracolsep{\fill}}lll@{}}
  \textbf{Course ID:} \courseid &
  \textbf{Instructor:} \instructorname &
  \textbf{Date:} \classdate
\end{tabular*}
\end{tcolorbox}%
}

\begin{document}
\pagestyle{empty} % Remove this line if page numbers are desired.

\makeclassnotesheader

\section{Convex functions}

Many inequalities in probability come from a single idea: a convex function
lies below its chords. This lecture makes that precise and draws out
Jensen's inequality and a few of its consequences.

\begin{statement}[Setting]
Throughout, $I \subseteq \mathbb{R}$ is an interval, $X$ is a random variable
taking values in $I$ with $\mathbb{E}\lvert X\rvert < \infty$, and
$\varphi \colon I \to \mathbb{R}$.
\end{statement}

\begin{definition}[Convex function]
The function $\varphi$ is \emph{convex} if for all $x, y \in I$ and
$\lambda \in [0,1]$,
\[
    \varphi\bigl(\lambda x + (1-\lambda) y\bigr) \le \lambda\varphi(x) + (1-\lambda)\varphi(y).
\]
It is \emph{strictly convex} if the inequality is strict whenever $x \ne y$
and $0 < \lambda < 1$.
\end{definition}

\begin{theorem}{Jensen's inequality}{jensen}
If $\varphi$ is convex and $\mathbb{E}\lvert\varphi(X)\rvert < \infty$, then
\[
    \varphi\bigl(\mathbb{E}[X]\bigr) \le \mathbb{E}\bigl[\varphi(X)\bigr].
\]
If $\varphi$ is strictly convex, equality holds only when $X$ is constant
almost surely.
\end{theorem}

\begin{proof}
Let $\mu = \mathbb{E}[X]$, which lies in $I$. A convex function has a
supporting line at every interior point: there is a slope $c$ with
\[
    \varphi(x) \ge \varphi(\mu) + c\,(x - \mu) \quad \text{for all } x \in I.
\]
Put $x = X$ and take expectations: $\mathbb{E}[\varphi(X)] \ge \varphi(\mu) + c\,(\mathbb{E}[X] - \mu) = \varphi(\mu)$.
\end{proof}

\begin{proposition}{Second-derivative test}{second-derivative}
If $\varphi$ is twice differentiable on the interior of $I$, then $\varphi$
is convex if and only if $\varphi''(x) \ge 0$ there, and strictly convex if
$\varphi'' > 0$.
\end{proposition}

\begin{question}
Which is larger, $\bigl(\mathbb{E}[X]\bigr)^2$ or $\mathbb{E}[X^2]$? And
$\log \mathbb{E}[X]$ or $\mathbb{E}[\log X]$, for $X > 0$?
\end{question}

\begin{answer}
Since $x \mapsto x^2$ is convex, $\bigl(\mathbb{E}[X]\bigr)^2 \le \mathbb{E}[X^2]$;
the difference is $\operatorname{Var}(X) \ge 0$. Since $\log$ is concave,
the inequality reverses: $\mathbb{E}[\log X] \le \log \mathbb{E}[X]$.
\end{answer}

\begin{example}[The AM--GM inequality]
Let $X$ take the values $a_1, \dots, a_n > 0$, each with probability $1/n$.
Jensen's inequality for the concave function $\log$ gives
\[
    \frac{1}{n}\sum_{i=1}^n \log a_i \le \log\Bigl(\frac{1}{n}\sum_{i=1}^n a_i\Bigr),
    \qquad\text{so}\qquad
    \sqrt[n]{a_1 a_2 \cdots a_n} \le \frac{a_1 + \cdots + a_n}{n}.
\]
\end{example}

\begin{remark}[Direction of the inequality]
A common slip is to apply Jensen's inequality in the wrong direction. Draw
the chord: for convex $\varphi$, the average of the values sits above the
value at the average.
\end{remark}

\end{document}